% =====================================================================
% Preamble additions required to compile methodology.tex
% Drop these \usepackage{} lines into the main paper preamble (alongside
% whatever the venue template already loads).  None of these clash with
% standard NeurIPS / ICML / ACL / ICLR / arXiv class files.
% =====================================================================

% --- math ---
\usepackage{amsmath,amssymb,amsfonts}
\usepackage{amsthm}
\newtheorem{remark}{Remark}        % \begin{remark}[Optional title] ... \end{remark}
\newtheorem{proposition}{Proposition}
\newtheorem{definition}{Definition}

% --- citations ---
% natbib gives you \citep{} and \citet{} which the section uses.
% If your venue template already loads natbib (most ML/NLP templates do),
% delete this line; the \citep{} / \citet{} calls will work unchanged.
\usepackage[authoryear,round]{natbib}

% --- algorithm / pseudocode ---
\usepackage{algorithm}
\usepackage{algpseudocode}
% (algpseudocode replaces the older algorithmic package and supports
%  \Statex, \Comment, \Require, \Ensure used in Algorithm 1.)

% --- tables ---
\usepackage{booktabs}     % \toprule / \midrule / \bottomrule

% --- figures (TikZ + libraries the architecture diagram uses) ---
\usepackage{tikz}
\usetikzlibrary{
  positioning,
  shapes.geometric,        % diamond gate
  arrows.meta,             % stealth arrowheads
  calc,                    % $(node)+(...)$ coordinate arithmetic
  fit,                     % optional grouping boxes
  backgrounds              % optional layered backgrounds
}

% --- code-style fonts (\texttt is sufficient; nothing extra needed) ---

% --- optional: tighter equation spacing if your template is loose ---
% \setlength{\abovedisplayskip}{4pt plus 1pt minus 1pt}
% \setlength{\belowdisplayskip}{4pt plus 1pt minus 1pt}

% =====================================================================
% Notes on cross-references the section assumes you have defined
% somewhere in the paper:
%
%   \label{sec:intro}        in Section 1 (Introduction)
%   \label{sec:experiments}  in Section 4 (Experiments)
%
% The methodology file itself defines every other label it references.
% =====================================================================
